\subsection{F01 Measurement Contract}
\label{app:measurement-contract}
The executed channel map is taken from the frozen SIM0001-JuliaPD-F01 research snapshot. The eight voltage buses are $\{2,5,6,10,19,22,29,39\}$. Voltage magnitude is reported in per unit with unit gain; voltage angle uses one common affine phase reference plus the archived relative offsets. The reference interpretation is observational: physical common drift and instrument clock drift are not separately identified. Terminal currents are converted to outward branch convention from the simulator\textquotesingle s into-terminal observable, then mapped with the row-specific polarity, gain and phase bias. Current bases use $100\,\mathrm{MVA}/(\sqrt{3}\,V_{\mathrm{terminal}})$; the common voltage base is 199.186 kV.\par
\begin{table*}[t]\centering\scriptsize\setlength{\tabcolsep}{3pt}
\caption{Eight PMU rows in the frozen F01 measurement map. V is a bus voltage in p.u.; I is an oriented branch-terminal phasor. Gains and phase biases are applied to current channels. BUS5 retains the historical 5--6, terminal-5 hypothesis and its producer CT identity is unverified. BUS39 current phase is partially validated.}
\label{tab:f01-measurement-contract}
\begin{tabular}{@{}rrrrllrrrll@{}}\toprule
PMU & V bus & Edge & Branch & Terminal & Side & Polarity & Gain & Bias ($^\circ$) & I unit & Provenance \\\midrule
2 & 2 & 2 & 2--3 & 3 & dst & 1 & 0.9581 & -1.173 & A & F01 archived terminal hypothesis; CT unverified \\
5 & 5 & 9 & 5--6 & 5 & src & 1 & 1.0361 & 0.312 & A & BUS5 5--6 historical hypothesis; CT unverified \\
6 & 6 & 11 & 6--7 & 6 & src & 1 & 1.0077 & 1.495 & A & F01 archived terminal hypothesis; CT unverified \\
10 & 10 & 17 & 10--11 & 11 & dst & -1 & 1.0041 & 2.224 & A & F01 archived terminal hypothesis; CT unverified \\
19 & 19 & 26 & 16--19 & 19 & dst & 1 & 1.0384 & 0.481 & A & F01 archived terminal hypothesis; CT unverified \\
22 & 22 & 34 & 21--22 & 22 & dst & 1 & 1.0189 & 0.094 & A & F01 archived terminal hypothesis; CT unverified \\
29 & 29 & 44 & 28--29 & 29 & dst & -1 & 1.0316 & -6.170 & A & F01 archived terminal hypothesis; CT unverified \\
39 & 39 & 1 & 1--39 & 39 & dst & -1 & 0.9298 & 30.053 & A & BUS39 partial phase validation \\
\bottomrule\end{tabular}
\end{table*}
Frequency is computed as $60+\dot\theta/(360)$ from the reported voltage angle, and ROCOF as $\ddot\theta/360$ using a causal 21-report quadratic estimate. The generator fit is robust and magnitude weighted; the predictor uses a fixed unweighted quadratic fit over the same history. Four warmup outputs are unavailable; the confirmatory evaluation uses a 10-s startup context and begins monitoring after 12 s. The 47 retained coordinates exclude the BUS39 current-angle coordinate. Noise multipliers 1 and 2 scale the standard deviation under the frozen stream. The stochastic model archives a cross-signal covariance field, but the executed report operator applies noise by signal family without that cross-family factor. These choices describe a simulated sensor configuration, not a verified eight-PMU CT/PT installation.
